\input{./._relpath-to-latexroot.ltx}
\documentclass[\latexroot/\projectname]{subfiles}
\input{\latexroot/@resources/subfile-setup.ltx}
\begin{document}

\whenintegrated{\label{microeconomic-literature}}\par\section{Wealth, Entrepreneurship, and Borrowing Constraints}
\whenintegrated{\label{sec:microlit}}

The original paper is motivated by a simple empirical fact: entrepreneurs are a small part of the population but hold a large share of wealth. I use those facts as background for the model reproduction, not as newly cleaned or newly reproduced SCF evidence in this project. The point for the current write-up is that the theory needs a mechanism that can connect business ownership, saving, and the upper tail of the wealth distribution.

Cagetti and De Nardi classify entrepreneurs using information on self-employment, business ownership, and active management. Table~\ref{tab:entrepreneurs-population} reports their benchmark definitions. The narrowest definition, self-employed business owners, is closest to the model idea of an active entrepreneur because it excludes households that only passively own a business.

\begin{table}[htbp]
  \centering
  \caption{Original Paper's Entrepreneurship Definitions and Wealth Shares}
  \whenintegrated{\label{tab:entrepreneurs-population}}
  \begin{tabular}{lcc}
    \hline\hline
    Definition & Percent in population & Share of total wealth \\
    \hline
    Business owners or self-employed & 16.7 & 52.9 \\
    All business owners & 13.3 & 48.8 \\
    Active business owners & 11.5 & 41.6 \\
    All self-employed & 11.1 & 39.0 \\
    Self-employed business owners & 7.6 & 33.0 \\
    \hline
  \end{tabular}
\end{table}

The model also starts from the observation that wealth is much more concentrated than labor earnings. Table~\ref{tab:wealth-distribution} gives the original paper's summary of the top of the U.S.\ wealth distribution. In this reproduction I treat these numbers as benchmark facts from the paper, not as outcomes of my own empirical data pipeline.

\begin{table}[htbp]
  \centering
  \caption{Benchmark Wealth Concentration Reported in the Original Paper}
  \whenintegrated{\label{tab:wealth-distribution}}
  \begin{tabular}{lcccc}
    \hline\hline
    & \multicolumn{4}{c}{Fraction of people, top} \\
    \cline{2-5}
    & $1\%$ & $5\%$ & $10\%$ & $20\%$ \\
    \hline
    Total net worth held & $30\%$ & $54\%$ & $67\%$ & $81\%$ \\
    \hline
  \end{tabular}
\end{table}

The literature before Cagetti and De Nardi had already made two points that are important for the reproduction. First, incomplete-markets models can generate precautionary saving but often have trouble matching the extreme upper tail of wealth. Second, entrepreneurship models show that liquidity and borrowing constraints can affect entry, investment, and saving. Examples include \cite{Evans1989-ws}, \cite{Quadrini2000-mp}, \cite{Gentry2004-rj}, and \cite{Buera2006-wj}. The contribution of the Cagetti and De Nardi model is to put occupational choice, entrepreneurial investment, borrowing constraints, and intergenerational wealth transmission inside one household problem.

For this project, the borrowing-constraint literature matters mainly because it tells me what to reproduce theoretically. The enforcement constraint is not just an exogenous debt limit. It is an incentive condition: the entrepreneur must prefer to repay rather than default and work as a wage earner. That is why the notebook gives special attention to the entrepreneur branch and to the comparison between the value of honoring debt and the value of default.

\smartbib

\end{document}
